\chapter{2013:Martin Karplus and Michael Levitt and Arieh Warshel}

\label{chap:chemistry2013}

The Nobel Prize in Chemistry for the year 2013 was awarded jointly to Martin Karplus and Michael Levitt and Arieh Warshel.

\textbf{Motivation:}

for the development of multiscale models for complex chemical systems

\section{Martin Karplus}

\subsection*{Early Life and Education}

Martin Karplus was born in 1930 in Vienna, Austria.

\subsection*{Career and Major Research}

Karplus earned a bachelor's degree from Harvard University in 1950 and a doctorate from the California Institute of Technology in 1953. He became a professor at Harvard University, where he pioneered the use of molecular dynamics to simulate the motion of macromolecules and co-developed methods that combine quantum and classical mechanics.

\subsection*{Nobel Prize-winning Work}

The work for which Martin Karplus was awarded the Nobel Prize in Chemistry in 2013 was described by the Nobel Committee as: "for the development of multiscale models for complex chemical systems".

Karplus showed how computer simulation could follow the movement of every atom in a large molecule, and his work underlies the molecular dynamics programs used across chemistry and biology.

\subsection*{Significance and Impact}

His simulations made it possible to watch how proteins move and fold and how they interact with other molecules, connecting theory with experiment at the atomic level.

\subsection*{Later Career and Honors}

Karplus was the Theodore William Richards Professor of Chemistry at Harvard University, where he worked for over five decades. He died on 2024-12-28 in Cambridge, Massachusetts, USA.

\subsection*{Legacy}

The computational methods he developed, including the CHARMM program, are standard tools for studying biomolecules and are used in drug design.

\subsection*{Selected Publications}

"CHARMM: A Program for Macromolecular Energy, Minimization, and Dynamics Calculations" (Journal of Computational Chemistry, 1983)

\clearpage

\section{Michael Levitt}

\subsection*{Early Life and Education}

Michael Levitt was born in 1947 in Pretoria, South Africa.

\subsection*{Career and Major Research}

Levitt studied at King's College London and received his doctorate from the University of Cambridge in 1971. He became a professor at Stanford University. He carried out some of the first molecular dynamics simulations of proteins and, with Arieh Warshel, developed a combined quantum and classical description of enzyme reactions.

\subsection*{Nobel Prize-winning Work}

The work for which Michael Levitt was awarded the Nobel Prize in Chemistry in 2013 was described by the Nobel Committee as: "for the development of multiscale models for complex chemical systems".

Levitt developed simplified representations of protein structure that made simulation of large systems feasible, and helped establish the multiscale approach that treats different parts of a molecule with different levels of theory.

\subsection*{Significance and Impact}

His simulations of protein dynamics and folding were among the first to show that computers could describe the behavior of biological macromolecules realistically.

\subsection*{Later Career and Honors}

Levitt continues his research on computational structural biology at Stanford University.

\subsection*{Legacy}

Levitt's methods are part of the foundation of computational structural biology, used for predicting protein structure and simulating biomolecular systems.

\subsection*{Selected Publications}

"Theoretical Studies of Enzymic Reactions: Dielectric, Electrostatic and Steric Stabilization of the Carbonium Ion in the Reaction of Lysozyme" (Journal of Molecular Biology, 1976)

\clearpage

\section{Arieh Warshel}

\subsection*{Early Life and Education}

Arieh Warshel was born in 1940 in Sde Nahum, Israel.

\subsection*{Career and Major Research}

Warshel earned bachelor's and master's degrees at the Technion--Israel Institute of Technology and a doctorate from the Weizmann Institute of Science in 1969. He became a professor at the University of Southern California, where he pioneered methods that combine quantum mechanics for reacting bonds with classical mechanics for the surrounding atoms.

\subsection*{Nobel Prize-winning Work}

The work for which Arieh Warshel was awarded the Nobel Prize in Chemistry in 2013 was described by the Nobel Committee as: "for the development of multiscale models for complex chemical systems".

Warshel developed the QM/MM approach, in which the part of a molecule where a chemical reaction occurs is treated quantum mechanically while the rest is treated classically, allowing enzyme reactions to be simulated.

\subsection*{Significance and Impact}

His calculations explained how enzymes accelerate reactions by stabilizing transition states, providing a quantitative link between enzyme structure and catalytic power.

\subsection*{Later Career and Honors}

Warshel continues his research on enzyme mechanisms and computational chemistry at the University of Southern California.

\subsection*{Legacy}

The QM/MM method is now used routinely to model reactions in enzymes, materials, and solution, and it is a cornerstone of modern computational chemistry.

\subsection*{Selected Publications}

"Theoretical Studies of Enzymic Reactions: Dielectric, Electrostatic and Steric Stabilization of the Carbonium Ion in the Reaction of Lysozyme" (Journal of Molecular Biology, 1976)

\clearpage

\section*{Chapter Summary}

The 2013 prize recognized the development of multiscale models that combine quantum and classical mechanics to simulate complex chemical systems. These methods allow chemists to compute how reactions proceed in enzymes and large molecules, a foundation of modern computational chemistry.
